\documentclass[12pt,a4paper]{exam}
\usepackage[utf8]{inputenc}
\usepackage{amsmath,amssymb,amsfonts}
\usepackage{geometry}
\usepackage{enumitem}
\usepackage{graphicx}
\usepackage{hyperref}
\usepackage{xcolor}
\geometry{a4paper, top=2cm, bottom=2cm, left=2.5cm, right=2.5cm}
\hypersetup{colorlinks=true, linkcolor=blue, urlcolor=blue}
\renewcommand{\thequestion}{\textbf{\arabic{question}}}
\renewcommand{\questionlabel}{\thequestion.}
\pointname{ marks}
\pointsdroppedatright
\bracketedpoints
\setlength{\parindent}{0pt}
\setlength{\emergencystretch}{3em}
\allowdisplaybreaks
\newsavebox{\schmathbox}
\newcommand{\fitmath}[1]{%
  \sbox{\schmathbox}{$\displaystyle #1$}%
  \ifdim\wd\schmathbox>\linewidth
    \resizebox{\linewidth}{!}{\usebox{\schmathbox}}%
  \else
    \usebox{\schmathbox}%
  \fi}

\begin{document}

\thispagestyle{empty}
\begin{center}
  \textbf{\LARGE Diag Solutions} \\[0.3cm]
  \rule{\textwidth}{0.5pt}
\end{center}

\vspace{0.3cm}

\begin{questions}
\question \textbf{1.} The angle of elevation of the top of a chimney from the top of a tower is $30^\circ$ and the angle of depression of the foot of the chimney from the top of the tower is $60^\circ$. If the tower is $40$ m high, find the height of the chimney.

\textbf{Answer:}
\begin{center}
\fitmath{160/3 \text{ m}}
\end{center}

\textbf{Solution:}
\begin{enumerate}
  \item Let $h$ be the height of the chimney above the tower level and $40$ m be the tower height.
  \item Distance $d = 40 \cot(60^\circ) = \frac{40}{\sqrt{3}}$.
  \item Height $h = d \tan(30^\circ) = \frac{40}{\sqrt{3}} \times \frac{1}{\sqrt{3}} = \frac{40}{3}$.
  \item Total height = $40 + \frac{40}{3} = \frac{160}{3}$ m.
\end{enumerate}
\textbf{Derivation:}
\begin{center}
\fitmath{H = 40 + \frac{40}{\sqrt{3}} \cdot \frac{1}{\sqrt{3}} = 40 + \frac{40}{3} = \frac{160}{3}}
\end{center}
\hfill \textit{Some Applications of Trigonometry — Application with two points}
\vspace{0.3cm}

\question \textbf{2.} A straight highway leads to the foot of a tower. A man standing at the top of the tower observes a car at an angle of depression of $30^\circ$, which is approaching the foot of the tower with a uniform speed. Six seconds later, the angle of depression of the car is found to be $60^\circ$. Find the time taken by the car to reach the foot of the tower from this point.

\textbf{Answer:}
\begin{center}
\fitmath{3 \text{seconds}}
\end{center}

\textbf{Solution:}
\begin{enumerate}
  \item Let $h$ be the height of the tower and $v$ be the uniform speed of the car.
  \item Let $A$ be the top of the tower and $BC$ be the distance covered by the car in 6 seconds.
  \item In $\triangle ABC$ (with angle $60^\circ$ at the car), $\tan 60^\circ = h/CD$.
  \item In $\triangle ABD$ (with angle $30^\circ$ at the car), $\tan 30^\circ = h/BD$.
  \item Relate the distances using the speed and time formula $d = v \times t$.
  \item Solve for the time $t$ taken to travel the remaining distance $CD$.
\end{enumerate}
\textbf{Derivation:}
Let  h  be height.  $\tan 60^\circ = h/x \implies x = h/\sqrt{3}$ .  $\tan 30^\circ = h/(x+d) \implies x+d = h\sqrt{3}$ . Substituting  x :  d = h$\sqrt{3} - h/\sqrt{3} = 2h/\sqrt{3}$ . Since distance  d  is covered in 6s, speed  v = d/6 = h/(3$\sqrt{3})$ . Time for distance  x  is  t = x/v = (h/$\sqrt{3}) / (h/3\sqrt{3}) = 3$  seconds.
\hfill \textit{Some Applications of Trigonometry — Angle of depression}
\vspace{0.3cm}

\question \textbf{3.} The angle of elevation of the top of a tower from a point on the ground is $30^\circ$. On walking $40$ metres towards the tower, the angle of elevation changes to $60^\circ$. Find the height of the tower and the distance of the initial point from the tower.

\textbf{Answer:}
\begin{center}
\fitmath{\text{Height} = 20\sqrt{3} m, \text{Distance} = 20\sqrt{3} + 40 m}
\end{center}

\textbf{Solution:}
\begin{enumerate}
  \item Let the height of the tower be $h$ and the distance from the second point be $x$.
  \item Set up the equation for the first point: $\tan 30^\circ = h / (x + 40)$.
  \item Set up the equation for the second point: $\tan 60^\circ = h / x$.
  \item Express $h$ in terms of $x$ from the second equation.
  \item Substitute into the first equation to solve for $x$ and then $h$.
\end{enumerate}
\textbf{Derivation:}
1)  h/x = $\sqrt{3} \implies h = x\sqrt{3}$ . 2)  h/(x+40) = 1/$\sqrt{3} \implies x\sqrt{3} / (x+40) = 1/\sqrt{3} \implies 3x = x+40 \implies 2x = 40 \implies x = 20$ . Height  h = 20$\sqrt{3}  m.$ Initial distance  = 20 + 40 = 60  m.
\hfill \textit{Some Applications of Trigonometry — Angle of elevation}
\vspace{0.3cm}

\question \textbf{4.} A vertical tower stands on a horizontal plane and is surmounted by a vertical flagstaff of height $6$ metres. At a point on the plane, the angle of elevation of the bottom and top of the flagstaff are $30^\circ$ and $60^\circ$ respectively. Find the height of the tower.

\textbf{Answer:}
\begin{center}
\fitmath{3 \text{metres}}
\end{center}

\textbf{Solution:}
\begin{enumerate}
  \item Let $h$ be the height of the tower and $d$ be the distance from the point to the base.
  \item Write the equation for the bottom of the flagstaff: $\tan 30^\circ = h/d$.
  \item Write the equation for the top of the flagstaff: $\tan 60^\circ = (h+6)/d$.
  \item Solve the system of equations for $h$.
\end{enumerate}
\textbf{Derivation:}
 d = h/$\tan 30^\circ = h\sqrt{3}$ . Substituting into second equation:  $\sqrt{3} = (h+6) / (h\sqrt{3}) \implies 3h = h+6 \implies 2h = 6 \implies h = 3  m.$
\hfill \textit{Some Applications of Trigonometry — Height of tower with flagstaff}
\vspace{0.3cm}

\question \textbf{5.} From the top of a $7$ m high building, the angle of elevation of the top of a cable tower is $60^\circ$ and the angle of depression of its foot is $45^\circ$. Determine the height of the tower.

\textbf{Answer:}
\begin{center}
\fitmath{7(1 + \sqrt{3}) \text{metres}}
\end{center}

\textbf{Solution:}
\begin{enumerate}
  \item Let the building be $AB = 7$ m and the tower be $CD$.
  \item Draw a horizontal line from $A$ to meet $CD$ at $E$. Thus $AE = BD$.
  \item In $\triangle ABD$, since the angle of depression is $45^\circ$, $BD = AB = 7$ m.
  \item In $\triangle ACE$, $\tan 60^\circ = CE / AE$.
  \item Calculate $CE$ and add $ED = AB = 7$ m to get the total height.
\end{enumerate}
\textbf{Derivation:}
 AE = BD = 7 . In  $\triangle$ ACE ,  CE = AE $\tan 60^\circ = 7\sqrt{3}$ . Total height  CD = CE + ED = 7$\sqrt{3} + 7 = 7(\sqrt{3} + 1)  m.$
\hfill \textit{Some Applications of Trigonometry — Combined elevation and depression}
\vspace{0.3cm}

\question \textbf{6.} A tangent $PQ$ at a point $P$ of a circle of radius $5 \text{ cm}$ meets a line through the centre $O$ at a point $Q$ so that $OQ = 13 \text{ cm}$. The length of $PQ$ is:

\textbf{Answer:}
(C)

\textbf{Solution:}
\begin{enumerate}
  \item The radius is perpendicular to the tangent at the point of contact.
  \item In the right-angled triangle $OPQ$, by Pythagoras theorem, $OQ^2 = OP^2 + PQ^2$.
  \item Substitute the values: $13^2 = 5^2 + PQ^2$.
  \item Solve for $PQ$.
\end{enumerate}
\textbf{Derivation:}
\begin{center}
\fitmath{PQ = \sqrt{OQ^2 - OP^2} = \sqrt{13^2 - 5^2} = \sqrt{169 - 25} = \sqrt{144} = 12 \text{ cm}}
\end{center}
\hfill \textit{Circles — Tangents to a circle}
\vspace{0.3cm}

\question \textbf{7.} If tangents $PA$ and $PB$ from a point $P$ to a circle with centre $O$ are inclined to each other at an angle of $80^\circ$, then $\angle POA$ is equal to:

\textbf{Answer:}
(A)

\textbf{Solution:}
\begin{enumerate}
  \item The line segment joining the centre to the external point bisects the angle between the tangents.
  \item Thus, $\angle APO = 80^\circ / 2 = 40^\circ$.
  \item In $\triangle OAP$, $\angle OAP = 90^\circ$ (tangent is perpendicular to radius).
  \item In $\triangle OAP$, $\angle POA = 180^\circ - 90^\circ - 40^\circ$.
\end{enumerate}
\textbf{Derivation:}
\begin{center}
\fitmath{\angle POA = 180^\circ - 90^\circ - 40^\circ = 50^\circ}
\end{center}
\hfill \textit{Circles — Properties of tangents}
\vspace{0.3cm}

\question \textbf{8.} From a point $Q$, the length of the tangent to a circle is $24 \text{ cm}$ and the distance of $Q$ from the centre is $25 \text{ cm}$. The radius of the circle is:

\textbf{Answer:}
(D)

\textbf{Solution:}
\begin{enumerate}
  \item Let $r$ be the radius. The triangle formed by the radius, tangent, and line to the centre is right-angled.
  \item $r^2 + 24^2 = 25^2$.
  \item $r^2 = 625 - 576 = 49$.
\end{enumerate}
\textbf{Derivation:}
\begin{center}
\fitmath{r = \sqrt{49} = 7 \text{ cm}}
\end{center}
\hfill \textit{Circles — Tangents to a circle}
\vspace{0.3cm}

\question \textbf{9.} Two concentric circles are of radii $5 \text{ cm}$ and $3 \text{ cm}$. The length of the chord of the larger circle which touches the smaller circle is:

\textbf{Answer:}
(B)

\textbf{Solution:}
\begin{enumerate}
  \item Let the chord be $AB$ and the point of contact be $M$. $OM \perp AB$.
  \item In $\triangle OMA$, $OA=5$ (hypotenuse) and $OM=3$ (radius).
  \item $AM = \sqrt{5^2 - 3^2} = 4$.
  \item Since the perpendicular from the centre bisects the chord, $AB = 2 \times AM$.
\end{enumerate}
\textbf{Derivation:}
\begin{center}
\fitmath{AB = 2 \times 4 = 8 \text{ cm}}
\end{center}
\hfill \textit{Circles — Concentric circles}
\vspace{0.3cm}

\question \textbf{10.} The number of tangents that can be drawn from a point inside a circle is:

\textbf{Answer:}
(C)

\textbf{Solution:}
\begin{enumerate}
  \item A tangent touches the circle at exactly one point.
  \item A point inside the circle lies in the interior region.
  \item Any line passing through an interior point will cut the circle at two points (secant).
  \item Therefore, no tangent can be drawn.
\end{enumerate}
\textbf{Derivation:}
Number of tangents = 0
\hfill \textit{Circles — Tangents from a point}
\vspace{0.3cm}

\question \textbf{11.} Assertion (A): The length of a tangent drawn from an external point to a circle is always equal to the radius of the circle.
Reason (R): The tangent at any point of a circle is perpendicular to the radius through the point of contact.

\textbf{Answer:}
(D)

\textbf{Solution:}
\begin{enumerate}
  \item Check Assertion: The length of a tangent from an external point is not necessarily equal to the radius; they are independent.
  \item Check Reason: This is a standard theorem (Theorem 10.1 of NCERT).
  \item Conclusion: Assertion is false, Reason is true.
\end{enumerate}
\textbf{Derivation:}
\begin{center}
\fitmath{A: \text{False}, R: \text{True}}
\end{center}
\hfill \textit{Circles — Tangents to a circle}
\vspace{0.3cm}

\question \textbf{12.} Assertion (A): If the distance between the centres of two circles of radii $r_1$ and $r_2$ is $d$, then they touch each other externally if $d = r_1 + r_2$.
Reason (R): The distance between the centres of two circles is the sum of their radii when they touch internally.

\textbf{Answer:}
(C)

\textbf{Solution:}
\begin{enumerate}
  \item Check Assertion: If two circles touch externally, the distance between their centres is indeed the sum of their radii. True.
  \item Check Reason: If two circles touch internally, the distance between their centres is $|r_1 - r_2|$. False.
\end{enumerate}
\textbf{Derivation:}
\begin{center}
\fitmath{A: \text{True}, R: \text{False}}
\end{center}
\hfill \textit{Circles — Relative positions of circles}
\vspace{0.3cm}

\question \textbf{13.} Assertion (A): The number of tangents that can be drawn from a point inside a circle is zero.
Reason (R): A tangent to a circle intersects the circle at only one point.

\textbf{Answer:}
(A)

\textbf{Solution:}
\begin{enumerate}
  \item Check Assertion: Points inside the circle cannot have tangents drawn to the circle. True.
  \item Check Reason: By definition, a tangent touches the circle at exactly one point. True.
  \item Conclusion: Reason explains why no tangent can be drawn from an internal point.
\end{enumerate}
\textbf{Derivation:}
A: True, R: True, R explains A
\hfill \textit{Circles — Tangents to a circle}
\vspace{0.3cm}

\question \textbf{14.} Assertion (A): The length of the tangent from a point $P$ at a distance of $10 \text{ cm}$ from the centre of a circle of radius $6 \text{ cm}$ is $8 \text{ cm}$.
Reason (R): In a right triangle, the square of the hypotenuse equals the sum of the squares of the other two sides.

\textbf{Answer:}
(A)

\textbf{Solution:}
\begin{enumerate}
  \item Check Assertion: $L = \sqrt{d^2 - r^2} = \sqrt{10^2 - 6^2} = \sqrt{100 - 36} = \sqrt{64} = 8$. True.
  \item Check Reason: Pythagoras theorem is used to calculate the tangent length. True.
\end{enumerate}
\textbf{Derivation:}
\begin{center}
\fitmath{L = \sqrt{10^2 - 6^2} = 8}
\end{center}
\hfill \textit{Circles — Tangents to a circle}
\vspace{0.3cm}

\question \textbf{15.} Assertion (A): Two parallel tangents can be drawn to a circle at the ends of a diameter.
Reason (R): Tangents at the ends of a diameter are perpendicular to the diameter.

\textbf{Answer:}
(A)

\textbf{Solution:}
\begin{enumerate}
  \item Check Assertion: Tangents at ends of diameter are parallel. True.
  \item Check Reason: Since the radius is perpendicular to the tangent at the point of contact, and the diameter is a straight line, the tangents are perpendicular to the diameter and thus parallel to each other. True.
\end{enumerate}
\textbf{Derivation:}
90\textasciicircum{}$\circ + 90^\circ = 180^\circ ($Co-interior angles)
\hfill \textit{Circles — Tangents to a circle}
\vspace{0.3cm}

\question \textbf{16.} Two concentric circles are of radii $5 \text{ cm}$ and $3 \text{ cm}$. Find the length of the chord of the larger circle which touches the smaller circle.

\textbf{Answer:}
\begin{center}
\fitmath{8 \text{ cm}}
\end{center}

\textbf{Solution:}
\begin{enumerate}
  \item Let $O$ be the center, $AB$ be the chord of the larger circle tangent to the smaller circle at $P$.
  \item Apply Pythagoras theorem in $\triangle OPA$ where $OA=5, OP=3$ to find $AP$, then double it.
\end{enumerate}
\textbf{Derivation:}
\begin{center}
\fitmath{AP = \sqrt{OA^2 - OP^2} = \sqrt{5^2 - 3^2} = \sqrt{25-9} = 4 \text{ cm}. \text{ Chord } AB = 2 \times AP = 2 \times 4 = 8 \text{ cm}.}
\end{center}
\hfill \textit{Circles — Tangent to a circle}
\vspace{0.3cm}

\question \textbf{17.} If a tangent $PA$ drawn from an external point $P$ to a circle with center $O$ makes an angle of $30^\circ$ with the chord $AB$, where $AB$ is the diameter, find $\angle APB$.

\textbf{Answer:}
\begin{center}
\fitmath{60^\circ}
\end{center}

\textbf{Solution:}
\begin{enumerate}
  \item Use the property that the angle between tangent and chord equals the angle in the alternate segment.
  \item In $\triangle PAB$, use the angle sum property or circle properties.
\end{enumerate}
\textbf{Derivation:}
\begin{center}
\fitmath{\angle ABP = 90^\circ (\text{angle in a semicircle}). \text{ Given } \angle PAB = 30^\circ. \text{ In } \triangle PAB, \angle APB = 180^\circ - (90^\circ + 30^\circ) = 60^\circ.}
\end{center}
\hfill \textit{Circles — Tangents and chords}
\vspace{0.3cm}

\question \textbf{18.} In the given figure, if $\angle OAB = 40^\circ$, find $\angle ACB$ where $O$ is the center of the circle.

\textbf{Answer:}
\begin{center}
\fitmath{50^\circ}
\end{center}

\textbf{Solution:}
\begin{enumerate}
  \item Identify that $\triangle OAB$ is isosceles with $OA=OB$.
  \item Find $\angle AOB$ and use the theorem that the angle at the center is double the angle at the circumference.
\end{enumerate}
\textbf{Derivation:}
\begin{center}
\fitmath{\angle OBA = \angle OAB = 40^\circ. \angle AOB = 180^\circ - (40^\circ + 40^\circ) = 100^\circ. \angle ACB = \frac{1}{2} \angle AOB = 50^\circ.}
\end{center}
\hfill \textit{Circles — Angle subtended by arc}
\vspace{0.3cm}

\question \textbf{19.} A quadrilateral $\text{ABCD}$ is drawn to circumscribe a circle. If $AB = 6 \text{ cm}$, $BC = 7 \text{ cm}$ and $CD = 4 \text{ cm}$, find the length of $AD$.

\textbf{Answer:}
\begin{center}
\fitmath{3 \text{ cm}}
\end{center}

\textbf{Solution:}
\begin{enumerate}
  \item Use the property that sums of opposite sides of a circumscribed quadrilateral are equal.
  \item Substitute the given values into $AB + CD = BC + AD$.
\end{enumerate}
\textbf{Derivation:}
\begin{center}
\fitmath{\begin{aligned}
AB + CD = BC + AD \\
\implies 6 + 4 = 7 + AD \\
\implies 10 = 7 + AD \\
\implies AD = 3 \text{ cm}.
\end{aligned}}
\end{center}
\hfill \textit{Circles — Tangents from external point}
\vspace{0.3cm}

\question \textbf{20.} Two tangents $TP$ and $TQ$ are drawn to a circle with center $O$ from an external point $T$. If $\angle PTQ = 70^\circ$, find $\angle POQ$.

\textbf{Answer:}
\begin{center}
\fitmath{110^\circ}
\end{center}

\textbf{Solution:}
\begin{enumerate}
  \item Use the property that the opposite angles of a quadrilateral formed by two tangents and two radii are supplementary.
  \item Subtract the given angle from $180^\circ$.
\end{enumerate}
\textbf{Derivation:}
\begin{center}
\fitmath{\angle OPT = \angle OQT = 90^\circ. \text{ In quadrilateral } \text{OPTQ}, \angle POQ + \angle PTQ = 180^\circ. \angle POQ = 180^\circ - 70^\circ = 110^\circ.}
\end{center}
\hfill \textit{Circles — Tangent properties}
\vspace{0.3cm}

\question \textbf{21.} A quadrilateral $\text{ABCD}$ is drawn to circumscribe a circle. Prove that $AB + CD = AD + BC$.

\textbf{Answer:}
\begin{center}
\fitmath{AB + CD = AD + BC}
\end{center}

\textbf{Solution:}
\begin{enumerate}
  \item Use the theorem that tangents drawn from an external point to a circle are equal in length.
  \item Identify pairs of equal tangents from vertices A, B, C, and D.
  \item Add the four equations to group the sides of the quadrilateral.
  \item Substitute the grouped segments to show the sum of opposite sides is equal.
\end{enumerate}
\textbf{Derivation:}
Let points of contact be P, Q, R, S on AB, BC, CD, DA respectively. AP=AS, BP=BQ, CQ=CR, DR=DS. Adding them: (AP+BP)+(CR+DR) = (AS+DS)+(BQ+CQ) => AB+CD = AD+BC.
\hfill \textit{Circles — Tangents to a circle}
\vspace{0.3cm}

\question \textbf{22.} Two concentric circles are of radii $5\text{ cm}$ and $3\text{ cm}$. Find the length of the chord of the larger circle which touches the smaller circle.

\textbf{Answer:}
\begin{center}
\fitmath{8\text{ cm}}
\end{center}

\textbf{Solution:}
\begin{enumerate}
  \item Draw the radius to the point of contact, which is perpendicular to the tangent.
  \item Form a right-angled triangle with the radius of the smaller circle, the radius of the larger circle, and half the chord.
  \item Apply Pythagoras theorem to find the length of half the chord.
  \item Multiply the result by 2 to get the full chord length.
\end{enumerate}
\textbf{Derivation:}
r\_1=5, r\_2=3. Let half-chord be x. x\textasciicircum{}2 + 3\textasciicircum{}2 = 5\textasciicircum{}2 => x\textasciicircum{}2 = 25 - 9 = 16 => x = 4. Chord length = 2x = 8$\text{ cm}.$
\hfill \textit{Circles — Tangents to a circle}
\vspace{0.3cm}

\question \textbf{23.} If a tangent $PA$ and $PB$ from a point $P$ to a circle with centre $O$ are inclined to each other at an angle of $80^\circ$, then find $\angle POA$.

\textbf{Answer:}
\begin{center}
\fitmath{50^\circ}
\end{center}

\textbf{Solution:}
\begin{enumerate}
  \item Identify that the line segment $OP$ bisects the angle between the tangents.
  \item Calculate the angle $\angle APO$ as half of $\angle APB$.
  \item Use the property that the radius is perpendicular to the tangent at the point of contact.
  \item Solve for $\angle POA$ in the right-angled triangle $\triangle OAP$.
\end{enumerate}
\textbf{Derivation:}
\begin{center}
\fitmath{\begin{aligned}
\angle APB = 80^\circ \\
\implies \angle APO = 40^\circ. \angle OAP = 90^\circ. \angle POA = 180^\circ - (90^\circ + 40^\circ) = 50^\circ.
\end{aligned}}
\end{center}
\hfill \textit{Circles — Tangents to a circle}
\vspace{0.3cm}

\question \textbf{24.} Prove that the lengths of tangents drawn from an external point to a circle are equal.

\textbf{Answer:}
\begin{center}
\fitmath{\text{Proof}}
\end{center}

\textbf{Solution:}
\begin{enumerate}
  \item Construct a circle with center O and external point P.
  \item Draw radii OA and OB to the points of contact A and B.
  \item Consider the two triangles $\triangle OAP$ and $\triangle OBP$.
  \item Prove congruence using RHS criterion and conclude PA = PB by CPCT.
\end{enumerate}
\textbf{Derivation:}
In $\triangle$ OAP and $\triangle$ OBP: OA=OB (radii), OP=OP (common), $\angle$ OAP=$\angle$ OBP=90\textasciicircum{}$\circ.$ By RHS, $\triangle$ OAP $\cong \triangle$ OBP. Thus, PA=PB.
\hfill \textit{Circles — Tangents to a circle}
\vspace{0.3cm}

\question \textbf{25.} A tangent $PQ$ at a point $P$ of a circle of radius $5\text{ cm}$ meets a line through the centre $O$ at a point $Q$ so that $OQ = 12\text{ cm}$. Find the length $PQ$.

\textbf{Answer:}
\begin{center}
\fitmath{\sqrt{119}\text{ cm}}
\end{center}

\textbf{Solution:}
\begin{enumerate}
  \item Recognize that the tangent $PQ$ is perpendicular to the radius $OP$.
  \item Identify $\triangle OPQ$ as a right-angled triangle with $\angle OPQ = 90^\circ$.
  \item Apply the Pythagoras theorem: $OQ^2 = OP^2 + PQ^2$.
  \item Substitute the given values and solve for $PQ$.
\end{enumerate}
\textbf{Derivation:}
\begin{center}
\fitmath{\begin{aligned}
OP=5, OQ=12. 12^2 = 5^2 + PQ^2 \\
\implies 144 = 25 + PQ^2 \\
\implies PQ^2 = 119 \\
\implies PQ = \sqrt{119}\text{ cm}.
\end{aligned}}
\end{center}
\hfill \textit{Circles — Tangents to a circle}
\vspace{0.3cm}

\question \textbf{26.} Prove that the lengths of tangents drawn from an external point to a circle are equal.

\textbf{Answer:}
\begin{center}
\fitmath{PA = PB}
\end{center}

\textbf{Solution:}
\begin{enumerate}
  \item Draw a circle with center O and an external point P.
  \item Draw tangents PA and PB touching the circle at A and B respectively.
  \item Join OA, OB, and OP.
  \item Consider triangles OAP and OBP.
  \item Show that the triangles are congruent by RHS congruence criterion.
  \item \(
\text{Conclude that PA} = \text{PB by CPCT}.
\)
\end{enumerate}
\textbf{Derivation:}
In  $\triangle$ OAP  and  $\triangle$ OBP :  OA = OB  (radii of the same circle).  $\angle$ OAP = $\angle$ OBP = 90\textasciicircum{}$\circ  ($tangent is perpendicular to radius at point of contact).  OP = OP  (common side). Therefore,  $\triangle$ OAP $\cong \triangle$ OBP  by RHS criterion. Hence,  PA = PB  by CPCT.
\hfill \textit{Circles — Tangents to a circle}
\vspace{0.3cm}

\question \textbf{27.} $PQ$ is a chord of length $8\text{ cm}$ of a circle of radius $5\text{ cm}$. The tangents at $P$ and $Q$ intersect at a point $T$. Find the length of $TP$.

\textbf{Answer:}
\begin{center}
\fitmath{TP = 20/3 cm}
\end{center}

\textbf{Solution:}
\begin{enumerate}
  \item Let the center be O. Join OP and OQ.
  \item Let OT intersect PQ at R. Since OP=OQ, $\triangle OPQ$ is isosceles.
  \item OT is the perpendicular bisector of PQ, so PR = 4 cm.
  \item In $\triangle OPR$, $OR = \sqrt{OP^2 - PR^2} = \sqrt{5^2 - 4^2} = 3\text{ cm}$.
  \item In $\triangle OPT$, $\angle OPT = 90^\circ$. Using similar triangles $\triangle PRT \sim \triangle OPT$.
  \item Set up the ratio: $TP/OP = PR/OR$. Solve for TP.
\end{enumerate}
\textbf{Derivation:}
In  $\triangle$ OPT ,  PR $\perp$ OT . Thus,  $\triangle$ PRT $\sim \triangle$ OPT  (AA similarity). Therefore,  TP/OP = PR/OR .  TP/5 = 4/3 .  TP = 20/3$\text{ cm}$ .
\hfill \textit{Circles — Properties of tangents}
\vspace{0.3cm}

\question \textbf{28.} A quadrilateral $\text{ABCD}$ is drawn to circumscribe a circle. Prove that $AB + CD = AD + BC$.

\textbf{Answer:}
\begin{center}
\fitmath{AB + CD = AD + BC}
\end{center}

\textbf{Solution:}
\begin{enumerate}
  \item Use the property that tangents from an external point to a circle are equal.
  \item Let the points of contact be P, Q, R, S on sides AB, BC, CD, DA respectively.
  \item \(
AP = AS, BP = BQ, CQ = CR, DR = DS.
\)
  \item Sum the left and right sides of these equations.
  \item Group the terms to form the sides of the quadrilateral.
  \item Conclude the result.
\end{enumerate}
\textbf{Derivation:}
\begin{center}
\fitmath{AP=AS, BP=BQ, CR=CQ, DR=DS. AB+CD = (AP+PB) + (CR+RD) = (AS+BQ) + (CQ+DS) = (AS+DS) + (BQ+CQ) = AD+BC.}
\end{center}
\hfill \textit{Circles — Circumscribed quadrilaterals}
\vspace{0.3cm}

\question \textbf{29.} Two concentric circles are of radii $5\text{ cm}$ and $3\text{ cm}$. Find the length of the chord of the larger circle which touches the smaller circle.

\textbf{Answer:}
\begin{center}
\fitmath{8 cm}
\end{center}

\textbf{Solution:}
\begin{enumerate}
  \item Let O be the common center. Let AB be the chord of the larger circle touching the smaller circle at point P.
  \item Join OP. Since AB is a tangent to the smaller circle, OP $\perp$ AB.
  \item OP = 3 cm (radius of smaller circle). OA = 5 cm (radius of larger circle).
  \item In right $\triangle OPA$, use Pythagoras theorem to find AP.
  \item AP = $\sqrt{OA^2 - OP^2} = \sqrt{5^2 - 3^2} = 4\text{ cm}$.
  \item Since the perpendicular from the center bisects the chord, AB = 2 * AP.
\end{enumerate}
\textbf{Derivation:}
In  $\triangle$ OPA ,  OA\textasciicircum{}2 = OP\textasciicircum{}2 + AP\textasciicircum{}2 $\implies 5^2 = 3^2 +$ AP\textasciicircum{}2 $\implies 25 = 9 +$ AP\textasciicircum{}2 $\implies$ AP\textasciicircum{}2 = 16 $\implies$ AP = 4$\text{ cm}$ . Since  OP $\perp$ AB ,  AP = PB = 4$\text{ cm}$ . Therefore,  AB = 4 + 4 = 8$\text{ cm}$ .
\hfill \textit{Circles — Concentric circles}
\vspace{0.3cm}

\question \textbf{30.} If the circumference of a circle is $176\text{ cm}$, then its radius is:

\textbf{Answer:}
(C)

\textbf{Solution:}
\begin{enumerate}
  \item Use the formula for circumference: $C = 2\pi r$
  \item Substitute $C = 176$ and $\pi = \frac{22}{7}$
  \item Solve for $r$
\end{enumerate}
\textbf{Derivation:}
\begin{center}
\fitmath{\begin{aligned}
176 = 2 \times \frac{22}{7} \times r \\
\implies 176 = \frac{44}{7} \times r \\
\implies r = \frac{176 \times 7}{44} = 4 \times 7 = 28 \text{ cm}
\end{aligned}}
\end{center}
\hfill \textit{Areas Related to Circles — Circumference of a Circle}
\vspace{0.3cm}

\question \textbf{31.} The area of a circle whose diameter is $14\text{ cm}$ is:

\textbf{Answer:}
(B)

\textbf{Solution:}
\begin{enumerate}
  \item Find the radius: $r = \frac{d}{2} = \frac{14}{2} = 7\text{ cm}$
  \item Use the area formula: $A = \pi r^2$
  \item Calculate $A$
\end{enumerate}
\textbf{Derivation:}
\begin{center}
\fitmath{A = \frac{22}{7} \times 7^2 = \frac{22}{7} \times 49 = 22 \times 7 = 154 \text{ cm}^2}
\end{center}
\hfill \textit{Areas Related to Circles — Area of a Circle}
\vspace{0.3cm}

\question \textbf{32.} The area of a quadrant of a circle with radius $7\text{ cm}$ is:

\textbf{Answer:}
(A)

\textbf{Solution:}
\begin{enumerate}
  \item Use the formula for area of a quadrant: $A = \frac{1}{4} \pi r^2$
  \item Substitute $r = 7$
  \item Calculate the result
\end{enumerate}
\textbf{Derivation:}
\begin{center}
\fitmath{A = \frac{1}{4} \times \frac{22}{7} \times 7^2 = \frac{1}{4} \times 154 = 38.5 \text{ cm}^2}
\end{center}
\hfill \textit{Areas Related to Circles — Area of a Sector}
\vspace{0.3cm}

\question \textbf{33.} If the perimeter of a semi-circular protractor is $36\text{ cm}$, then its diameter is:

\textbf{Answer:}
(D)

\textbf{Solution:}
\begin{enumerate}
  \item The perimeter of a semi-circle is $\pi r + 2r = r(\pi + 2)$
  \item Set $r(\frac{22}{7} + 2) = 36$
  \item Solve for $r$ and find $d = 2r$
\end{enumerate}
\textbf{Derivation:}
\begin{center}
\fitmath{\begin{aligned}
r(\frac{22+14}{7}) = 36 \\
\implies r(\frac{36}{7}) = 36 \\
\implies r = 7 \\
\implies d = 2 \times 7 = 14 \text{ cm}
\end{aligned}}
\end{center}
\hfill \textit{Areas Related to Circles — Circumference of a Circle}
\vspace{0.3cm}

\question \textbf{34.} The area of a sector of a circle with radius $6\text{ cm}$ and central angle $60^\circ$ is:

\textbf{Answer:}
(C)

\textbf{Solution:}
\begin{enumerate}
  \item Use the formula: $A = \frac{\theta}{360^\circ} \times \pi r^2$
  \item Substitute $\theta = 60^\circ$ and $r = 6$
  \item Calculate
\end{enumerate}
\textbf{Derivation:}
\begin{center}
\fitmath{A = \frac{60}{360} \times \pi \times 6^2 = \frac{1}{6} \times \pi \times 36 = 6\pi \text{ cm}^2}
\end{center}
\hfill \textit{Areas Related to Circles — Area of a Sector}
\vspace{0.3cm}

\question \textbf{35.} Assertion (A): The area of a circle with radius $7$ cm is $154$ cm$^2$. Reason (R): The area of a circle is given by $\pi r^2$, where $r$ is the radius.

\textbf{Answer:}
(A)

\textbf{Solution:}
\begin{enumerate}
  \item Formula for area of a circle is $\pi r^2$.
  \item Substitute $r = 7$: $\frac{22}{7} \times 7 \times 7 = 22 \times 7 = 154$.
\end{enumerate}
\textbf{Derivation:}
\begin{center}
\fitmath{\text{Area} = \pi r^2 = \frac{22}{7} \times (7)^2 = 154 \text{ cm}^2}
\end{center}
\hfill \textit{Areas Related to Circles — Area of a circle}
\vspace{0.3cm}

\question \textbf{36.} Assertion (A): The circumference of a circle with diameter $14$ cm is $44$ cm. Reason (R): The circumference of a circle is given by $\pi d$, where $d$ is the diameter.

\textbf{Answer:}
(A)

\textbf{Solution:}
\begin{enumerate}
  \item Circumference formula is $\pi d$.
  \item Substitute $d = 14$: $\frac{22}{7} \times 14 = 22 \times 2 = 44$.
\end{enumerate}
\textbf{Derivation:}
\begin{center}
\fitmath{C = \pi d = \frac{22}{7} \times 14 = 44 \text{ cm}}
\end{center}
\hfill \textit{Areas Related to Circles — Circumference of a circle}
\vspace{0.3cm}

\question \textbf{37.} Assertion (A): The area of a semicircle with radius $r$ is $\frac{1}{2} \pi r^2$. Reason (R): The circumference of a semicircle is $\pi r$.

\textbf{Answer:}
(C)

\textbf{Solution:}
\begin{enumerate}
  \item Area of semicircle is indeed $\frac{1}{2} \pi r^2$ (True).
  \item Circumference of a semicircle includes the diameter, so it is $\pi r + 2r$ (False).
\end{enumerate}
\textbf{Derivation:}
\begin{center}
\fitmath{\text{Area} = \frac{1}{2} \pi r^2; \text{Perimeter} = \pi r + 2r}
\end{center}
\hfill \textit{Areas Related to Circles — Area of semicircle}
\vspace{0.3cm}

\question \textbf{38.} Assertion (A): If the perimeter and the area of a circle are numerically equal, then the radius of the circle is $2$ units. Reason (R): Circumference $2\pi r = \pi r^2$ implies $r=2$.

\textbf{Answer:}
(A)

\textbf{Solution:}
\begin{enumerate}
  \item Set $2\pi r = \pi r^2$.
  \item Divide by $\pi r$ (since $r \neq 0$), we get $2 = r$.
\end{enumerate}
\textbf{Derivation:}
\begin{center}
\fitmath{\begin{aligned}
2\pi r = \pi r^2 \\
\implies r = \frac{2\pi r}{\pi r} = 2
\end{aligned}}
\end{center}
\hfill \textit{Areas Related to Circles — Numerical equality of area and circumference}
\vspace{0.3cm}

\question \textbf{39.} Assertion (A): The area of a sector of a circle with radius $r$ and central angle $\theta$ is $\frac{\theta}{360} \times \pi r^2$. Reason (R): The length of the arc of a sector is $\frac{\theta}{360} \times \pi r^2$.

\textbf{Answer:}
(C)

\textbf{Solution:}
\begin{enumerate}
  \item Area of sector formula is correct (True).
  \item Arc length formula is $\frac{\theta}{360} \times 2\pi r$ (False).
\end{enumerate}
\textbf{Derivation:}
\begin{center}
\fitmath{\text{Area} = \frac{\theta}{360} \pi r^2; \text{Length} = \frac{\theta}{360} 2\pi r}
\end{center}
\hfill \textit{Areas Related to Circles — Sector of a circle}
\vspace{0.3cm}

\question \textbf{40.} Find the area of a sector of a circle with radius $6 \text{ cm}$ if the angle of the sector is $60^\circ$. (Use $\pi = 3.14$)

\textbf{Answer:}
\begin{center}
\fitmath{18.84 \text{ cm}^2}
\end{center}

\textbf{Solution:}
\begin{enumerate}
  \item Use the formula for the area of a sector: $A = \frac{\theta}{360^\circ} \times \pi r^2$.
  \item Substitute the given values $r=6$ and $\theta=60$ into the formula and calculate.
\end{enumerate}
\textbf{Derivation:}
\begin{center}
\fitmath{A = \frac{60}{360} \times 3.14 \times 6^2 = \frac{1}{6} \times 3.14 \times 36 = 3.14 \times 6 = 18.84 \text{ cm}^2}
\end{center}
\hfill \textit{Areas Related to Circles — Area of a Sector}
\vspace{0.3cm}

\end{questions}

\end{document}
